% Sección 12.2: Límites y Continuidad

\subsection*{Ejercicios 1--2}
\begin{enumerate}
    \subimport{ejercicio_01/}{ejercicio.tex}
    \subimport{ejercicio_02/}{ejercicio.tex}
\end{enumerate}

\subsection*{Ejercicios 3--22: Encuentre el límite, si existe, o demuestre que el límite no existe.}
\begin{enumerate}
    \setcounter{enumi}{2}
    \subimport{ejercicio_03/}{ejercicio.tex}
    \subimport{ejercicio_04/}{ejercicio.tex}
    \subimport{ejercicio_05/}{ejercicio.tex}
    \subimport{ejercicio_06/}{ejercicio.tex}
    \subimport{ejercicio_07/}{ejercicio.tex}
    \subimport{ejercicio_08/}{ejercicio.tex}
    \subimport{ejercicio_09/}{ejercicio.tex}
    \subimport{ejercicio_10/}{ejercicio.tex}
    \subimport{ejercicio_11/}{ejercicio.tex}
    \subimport{ejercicio_12/}{ejercicio.tex}
    \subimport{ejercicio_13/}{ejercicio.tex}
    \subimport{ejercicio_14/}{ejercicio.tex}
    \subimport{ejercicio_15/}{ejercicio.tex}
    \subimport{ejercicio_16/}{ejercicio.tex}
    \subimport{ejercicio_17/}{ejercicio.tex}
    \subimport{ejercicio_18/}{ejercicio.tex}
    \subimport{ejercicio_19/}{ejercicio.tex}
    \subimport{ejercicio_20/}{ejercicio.tex}
    \subimport{ejercicio_21/}{ejercicio.tex}
    \subimport{ejercicio_22/}{ejercicio.tex}
\end{enumerate}

\subsection*{Ejercicios 23--24: Encuentre el límite, si existe, determinando los valores correspondientes para funciones de tres variables.}
\begin{enumerate}
    \setcounter{enumi}{22}
    \subimport{ejercicio_23/}{ejercicio.tex}
    \subimport{ejercicio_24/}{ejercicio.tex}
\end{enumerate}

\subsection*{Ejercicios 25--28: Utilice coordenadas polares para encontrar el límite, si existe.}
\begin{enumerate}
    \setcounter{enumi}{24}
    \subimport{ejercicio_25/}{ejercicio.tex}
    \subimport{ejercicio_26/}{ejercicio.tex}
    \subimport{ejercicio_27/}{ejercicio.tex}
    \subimport{ejercicio_28/}{ejercicio.tex}
\end{enumerate}

\subsection*{Ejercicios 29--38: Determine la región de continuidad de la función dada.}
\begin{enumerate}
    \setcounter{enumi}{28}
    \subimport{ejercicio_29/}{ejercicio.tex}
    \subimport{ejercicio_30/}{ejercicio.tex}
    \subimport{ejercicio_31/}{ejercicio.tex}
    \subimport{ejercicio_32/}{ejercicio.tex}
    \subimport{ejercicio_33/}{ejercicio.tex}
    \subimport{ejercicio_34/}{ejercicio.tex}
    \subimport{ejercicio_35/}{ejercicio.tex}
    \subimport{ejercicio_36/}{ejercicio.tex}
    \subimport{ejercicio_37/}{ejercicio.tex}
    \subimport{ejercicio_38/}{ejercicio.tex}
\end{enumerate}

\subsection*{Ejercicios 39--40: Determine si la función es continua en el origen.}
\begin{enumerate}
    \setcounter{enumi}{38}
    \subimport{ejercicio_39/}{ejercicio.tex}
    \subimport{ejercicio_40/}{ejercicio.tex}
\end{enumerate}

\subsection*{Ejercicios 41--42: Utilice coordenadas polares para demostrar que la función tiene una discontinuidad removible en el origen y defina $f(0,0)$ de manera adecuada para que sea continua.}
\begin{enumerate}
    \setcounter{enumi}{40}
    \subimport{ejercicio_41/}{ejercicio.tex}
    \subimport{ejercicio_42/}{ejercicio.tex}
\end{enumerate}
